\section{Cost composition and the next research obligations}
\label{sec:research-agenda}

The results above improve two representation costs: extracting the
components of a supplied normal surface, and retaining a complete block
of raw singleton substitutions. The next theorem states precisely how
such local improvements could contribute to a complete recognition bound.
It also makes clear which quantities future work must control.

\subsection{A conditional theorem for the whole computation}

Let $n\ge2$ be the bit length of a diagram encoding. Standard encodings
whose size is polynomial in the crossing count give the same
quasi-polynomial complexity class. All subsequent sizes are measured in
bits, including integer coordinates, combinatorial incidence, word
circuits, retained history and proof references.

\begin{proposition}[Conditional cost composition]
\label{prop:conditional-cost-composition}
Suppose a deterministic algorithm constructs and evaluates a search tree
for an input of size $n$, subject to the following uniform bounds.
The transition and terminal rules are sound, and exhausting the tree is
complete for both unknot and knotted verdicts. The tree has depth at most
$h=h(n)$, and each node at depth $j<h$ has at most $b_j=b_j(n)\ge1$
children. Failed candidate attempts and repair iterations are included
in the tree or in its explicitly charged node work.

Every intermediate state has at most $s=s(n)\ge1$ bits. All work at
one node, including generation of its candidates and rejected proposals,
has bit cost polynomial in $n+s$, with one polynomial independent of the
input and node. Each node contributes at most polynomially many bits in
$n+s$ to a source-bound proof record. The assembled proof is independently
checkable in time polynomial in its full length and $n$.

Then there is an absolute constant $a$ such that discovery, proof output,
and independent verification together have cost at most
\begin{equation}\label{eq:global-composition-cost}
 \left((h+1)(n+s+2)\prod_{j=0}^{h-1}b_j\right)^a.
\end{equation}
In particular, for a fixed $c\ge1$, the sufficient condition
\begin{equation}\label{eq:global-log-budget}
 \log_2(n+s+2)+\log_2(h+1)
       +\sum_{j=0}^{h-1}\log_2 b_j
       =O\bigl((\log n)^c\bigr)
\end{equation}
gives a complete deterministic recognition algorithm of complexity
$2^{O((\log n)^c)}$.
\end{proposition}

\begin{proof}
The number of nodes at depth $j$ is at most
$\prod_{i=0}^{j-1}b_i$. Since every $b_i\ge1$, the total node count is
at most $V=(h+1)\prod_{i=0}^{h-1}b_i$. Summing the uniform node cost
and proof length gives $V(n+s+2)^{O(1)}$ for each. Applying the assumed
polynomial verifier to that full proof preserves a fixed polynomial in
$V$ and $n+s+2$, proving~\eqref{eq:global-composition-cost}. Taking
logarithms gives~\eqref{eq:global-log-budget}. Soundness and completeness,
which are separate hypotheses, turn this computation bound into a
recognition bound.
\end{proof}

For example, $s=n^{O(1)}$, $b_j=n^{O(1)}$ uniformly, and
$h=O(\log n)$ give $2^{O((\log n)^2)}$. Replacing the depth by
$O((\log n)^2)$ gives $2^{O((\log n)^3)}$. A deterministic chain with
polynomially many steps and polynomial state size needs no logarithmic-depth
hypothesis: its branching factors are all one. A small description of each
state does not establish the branching term; polynomial-time operations
do not establish a bound on the largest state. If the terminal rules
only certify unknottedness, the same analysis bounds a possibly
inconclusive positive search and does not supply the missing negative
verdict.

The analogous accounting for an iteration bound $L(g+1)^L$, as in
Lackenby's hierarchy analysis~\cite[Proposition 9.1]{lackenby2026}, is
\[
 \log_2(n+s+2)+\log_2(L+1)+L\log_2(g+1)
                  =O((\log n)^c).
\]
This is conditional on bounded candidate generation, repairs and replay
throughout the computation. Bounds for an essential boundary pattern
cannot simply be applied to arbitrary unsuccessful branches.

\subsection{Prioritized questions and concrete targets}

The following questions are ordered by their role in completing the
geometric and algebraic pipelines. Each asks for a specific constructive
result or obstruction, rather than an improvement to isolated timings.

\begin{enumerate}
\item \textbf{How can a useful surface be found?}
Can one generate a complete family of useful normal or hierarchy surfaces
after $2^{\operatorname{polylog} n}$ candidate attempts, with uniformly
bounded encoded coordinates? Component extraction now permits a candidate
to be disconnected. The missing statement concerns coverage: whenever
the required compression or hierarchy move exists, some generated
candidate must contain it. A first result should specify an actual
diagram or triangulation class and a parameter controlling this search,
including the cost of finding that parameter's witnessing structure.

\item \textbf{Can cutting be performed directly on compressed coordinates?}
Given a validated triangulation and a binary normal vector, construct a
compressed cell or handle description of the cut manifold, together with
its boundary pattern and a checkable identification with the geometric
cut. The target is polynomial cost in the input and output bit lengths,
without allocating one cell per normal disc. Weighted segment tracking
under crushing already has a concrete precedent in the work of
He, Sedgwick and Spreer~\cite[Section 4.1]{hesedgwickspreer2025}; a full
checked cut representation requires additional topology and attachment
data. State exactly which parts
require essentiality hypotheses and provide a defined outcome when a
candidate does not satisfy them.

\item \textbf{Which ordered attachment data admit small descriptions?}
Extend sparse component incidence to boundary order and transport of
named arcs, points and intervals. An acceptable representation must
support the actual gluing operations, not merely count incidences.
Determine whether a useful class has polynomial-size attachment maps,
and construct lower-bound examples outside it. The linear bound on
unmarked component vectors does not bound the number of possible
markings; relative offsets alone can distinguish many states.

\item \textbf{How is peripheral provenance maintained?}
Connect the input diagram to a finite exterior triangulation with checked
meridian and longitude data. Then carry that data through cutting,
crushing, retriangulation and any change of combinatorial coordinates.
The desired output is a composable source-bound certificate for the
peripheral identifications. Recognizing an essential disc in a supplied
torus-boundary manifold is insufficient until that manifold and its
markings are linked to the original diagram.

\item \textbf{How small can a selected-component witness be?}
The complete coalesced inventory is polynomial and has at most $v+2q$
types. Can a particular disc or marked component be certified with a
substantially smaller record than the whole weighted trace, while still
proving connectedness and source membership? Investigate persistent
anchors, locally checkable transport summaries, and streaming replay.
Any savings must include locating the component and verifying its boundary
class; conservation of the total vector is not a substitute.

\item \textbf{Can the geometric checker have a smaller trusted core?}
Separate combinatorial manifold validation, interval equivalence,
integer weight transport and surface classification into explicit proof
interfaces. Formalize the local transport and normal-anchor statements,
or implement an independently derived second checker. Measure proof
bytes and maximum replay state as well as arithmetic work. Resource
limits must preserve the distinction between an invalid certificate and
an unfinished computation, and malformed records must be charged by
their actual encoded length.

\item \textbf{Can persistent elimination be integrated without changing semantics?}
Add source-derived replay for immutable concatenations and generator
bindings, retaining exact relator-slot identities and the raw singleton
condition. Produce an end-to-end diagram experiment in which discovery
and independent replay both use this representation. Ordinary SLP export
provides a clear boundary to other word kernels. The mathematical target
is an integrated polynomial bound for each complete raw block, including
donor identification and proof construction, rather than the substitution
operation alone.

\item \textbf{What controls growth across normalization boundaries?}
A polynomial bound for one raw block does not automatically survive
repeated free reduction, cyclic normalization and reconstruction.
Seek an amortized potential controlling all allocated grammar vertices,
integer metadata and proof records across the entire sequence. A useful
theorem should bound the total composition, not repeatedly apply a fixed
polynomial whose degree compounds with the number of boundaries.
Counterexamples for particular normalization policies would identify
which kinds of sharing must be retained.

\item \textbf{When does donor planning admit a complete efficient policy?}
Computing all current raw singleton counts is polynomial, but choosing
an order that reaches a terminal presentation may involve many choices.
Identify an exchange property, a monotone dependency structure, or a
restricted-width class on which a specified policy is complete. Charge
failed exposure moves and alternative orders explicitly. A promising
first theorem would connect such an algebraic class to a recognizable
family of diagrams, rather than assuming a successful donor sequence
as additional input.

\item \textbf{Can hierarchy repairs be accelerated with a proved potential?}
Construct the actual correction operations and bound both their number
and the sizes of the descriptions they create. A logarithmic-depth
theorem needs a quantity that decreases under the accelerated procedure,
including unsuccessful choices and subsequent repairs. Determine whether
the potential should combine handle complexity, boundary-pattern
complexity and compressed attachment size. A terminating integer
potential whose initial value is exponential does not meet this target.

\item \textbf{How do the geometric and algebraic branching costs compose?}
Establish bounds for the sum of logarithmic branching factors in
\eqref{eq:global-log-budget}, allowing depth-dependent choices and shared
states. Test whether reductions in one representation preserve the
parameters that make the other efficient. A restricted theorem may be
valuable even without a universal result, provided its admission
conditions are recognizable and remain valid after every permitted
transition. Switching between individually efficient backends is not
itself a bound on the resulting search.

\item \textbf{Which evidence best tests the proposed global bounds?}
Build reproducible diagram families that vary coordinate bit length,
distinct component types, attachment complexity, raw elimination depth
and normalization frequency independently. Record candidate failures,
conversion, proof generation, proof bytes, replay and fallback. Include
cases designed to saturate proved bounds and cases that defeat favorable
heuristics. Such experiments should guide the next structural conjecture
and identify its weakest hypothesis; a successful finite corpus cannot
replace the required completeness and worst-case arguments.
\end{enumerate}

